\documentclass[11pt,a4paper]{scrreprt}

% ---------- Layout ----------
\usepackage[a4paper,top=2cm,bottom=2.1cm,left=2.2cm,right=2.2cm,headheight=15pt]{geometry}
\usepackage[utf8]{inputenc}
\usepackage[T1]{fontenc}
\usepackage{lmodern}

% ---------- Color ----------
\usepackage{xcolor}
\definecolor{deepnavy}{RGB}{25,52,84}
\definecolor{tealaccent}{RGB}{40,132,132}
\definecolor{warmgold}{RGB}{198,148,42}
\definecolor{softgray}{RGB}{244,244,240}
\definecolor{inkgray}{RGB}{70,70,70}

% ---------- Drawing ----------
\usepackage{tikz}
\usetikzlibrary{positioning,calc}
\usepackage[most]{tcolorbox}
\tcbuselibrary{skins}

% ---------- Text / math / tables ----------
\usepackage{amsmath}
\usepackage{amssymb}
\usepackage{booktabs}
\usepackage{array}
\usepackage{enumitem}
\usepackage{fontawesome5}

% ---------- Headers / footers ----------
\usepackage[automark]{scrlayer-scrpage}
\pagestyle{scrheadings}
\clearpairofpagestyles
\ihead{\small\color{inkgray}\textsc{Distance Learning Study Guide}}
\ohead{\small\color{inkgray}ENV\,101 -- Module Guide}
\ofoot{\small\color{inkgray}Page \thepage\ of \pageref{LastPage}}
\ifoot{\small\color{inkgray}Meridian Distance Learning Institute}
\usepackage{lastpage}

% ---------- KOMA heading styling ----------
\addtokomafont{disposition}{\color{deepnavy}}
\addtokomafont{chapter}{\Huge\bfseries}
\addtokomafont{section}{\Large\bfseries\color{tealaccent}}
\addtokomafont{subsection}{\large\bfseries}
\setkomafont{chapterentry}{\normalfont}
\RedeclareSectionCommand[
  beforeskip=0.4cm,
  afterskip=0.6cm
]{chapter}

\usepackage[hidelinks,colorlinks=true,linkcolor=deepnavy,urlcolor=tealaccent]{hyperref}

% ---------- Reusable box styles ----------
\newtcolorbox{objectivesbox}{
  enhanced,
  breakable,
  colback=softgray,
  colframe=deepnavy,
  boxrule=0.6pt,
  left=9pt,right=9pt,top=5pt,bottom=5pt,
  title={\faBullseye\ \ Learning Objectives},
  fontupper=\small,
  fonttitle=\bfseries\color{white},
  colbacktitle=deepnavy,
  arc=2pt
}

\newtcolorbox{infobox}[1]{
  enhanced,
  breakable,
  colback=white,
  colframe=tealaccent,
  boxrule=0.6pt,
  left=8pt,right=8pt,top=4pt,bottom=4pt,
  title=#1,
  fonttitle=\bfseries\color{white},
  colbacktitle=tealaccent,
  arc=2pt
}

\newtcolorbox{termsbox}{
  enhanced,
  breakable,
  colback=softgray,
  colframe=warmgold,
  boxrule=0.6pt,
  left=9pt,right=9pt,top=5pt,bottom=5pt,
  title={\faBook\ \ Key Terms},
  fontupper=\small,
  fonttitle=\bfseries\color{white},
  colbacktitle=warmgold,
  arc=2pt
}

\newtcolorbox{quizbox}{
  enhanced,
  breakable,
  colback=white,
  colframe=deepnavy,
  boxrule=0.8pt,
  left=9pt,right=9pt,top=5pt,bottom=5pt,
  title={\faQuestionCircle\ \ Self-Check Quiz},
  fontupper=\small,
  fonttitle=\bfseries\color{white},
  colbacktitle=deepnavy,
  arc=2pt
}

\newtcolorbox{activitybox}{
  enhanced,
  breakable,
  colback=softgray,
  colframe=tealaccent,
  boxrule=0.6pt,
  left=9pt,right=9pt,top=5pt,bottom=5pt,
  title={\faPen\ \ Applied Activity},
  fontupper=\small,
  fonttitle=\bfseries\color{white},
  colbacktitle=tealaccent,
  arc=2pt
}

\setlist[itemize,1]{leftmargin=1.4em,itemsep=2pt,topsep=2pt}
\setlist[enumerate,1]{leftmargin=1.6em,itemsep=3pt,topsep=2pt}

\begin{document}

% ================= TITLE PAGE =================
\begin{titlepage}
\begin{tikzpicture}[remember picture,overlay]
  \fill[deepnavy] (current page.north west) rectangle ([yshift=-4.2cm]current page.north east);
  \fill[warmgold] ([yshift=-4.2cm]current page.north west) rectangle ([yshift=-4.28cm]current page.north east);
\end{tikzpicture}

\vspace*{1.1cm}
\begin{flushleft}
\hspace{1.6cm}\textcolor{white}{\Large\textsc{Meridian Distance Learning Institute}}
\end{flushleft}

\vspace{3.4cm}
\begin{center}
{\Huge\bfseries\color{deepnavy} Foundations of\\[2pt] Environmental Science}\\[10pt]
{\Large\color{tealaccent} Self-Paced Study Guide}\\[4pt]
{\large\color{inkgray} Course Code ENV\,101 \quad$\vert$\quad Modules 1--2 of 8}
\end{center}

\vspace{1.6cm}
\begin{center}
\begin{tikzpicture}
\draw[deepnavy,line width=0.8pt] (0,0) -- (11.5,0);
\end{tikzpicture}
\end{center}

\vspace{1cm}
\begin{center}
\begin{minipage}{0.78\textwidth}
\begin{infobox}{\faGraduationCap\ \ Course At a Glance}
\begin{tabular}{@{}p{3.4cm}p{7.6cm}@{}}
\textbf{Instructor:} & Dr.\ Elena Whitfield, Ph.D.\ Environmental Systems \\
\textbf{Credit Load:} & 3 semester credits \\
\textbf{Format:} & Fully asynchronous; readings, recorded lectures, quizzes \\
\textbf{Estimated Pace:} & 3--4 hours per module, 8 modules total \\
\textbf{Support:} & Weekly virtual office hours, discussion board within 24h \\
\end{tabular}
\end{infobox}
\end{minipage}
\end{center}

\vspace{10pt}
\begin{center}
\begin{minipage}{0.78\textwidth}
\begin{infobox}{\faLightbulb\ \ How to Use This Guide}
\small
This guide mirrors the online course shell module by module. Every chapter opens with
\textbf{\faBullseye\ objectives}, lists an \textbf{\faClock\ estimated time} and reading, closes
practice with a \textbf{\faPen\ hands-on activity}, and checks understanding with a
\textbf{\faQuestionCircle\ self-check quiz} -- answers are in the Answer Key at the very end.
Do the activity before the quiz; skipping it is the single biggest predictor of a failed
self-check.
\end{infobox}
\end{minipage}
\end{center}

\vfill
\begin{center}
{\small\color{inkgray} This guide accompanies the ENV\,101 course shell and the required text\\
\emph{Systems of the Living Earth}, 4th ed.\ (Nimbus Academic Press).}\\[6pt]
{\small\color{inkgray} Revision 3 -- Spring Term}
\end{center}
\end{titlepage}

% ================= CHAPTER 1 =================
\chapter{Earth as an Interconnected System}

\begin{objectivesbox}
By the end of this chapter, you will be able to:
\begin{enumerate}[label=\arabic*.]
  \item Identify the four major spheres of the Earth system and describe how they interact.
  \item Explain the difference between an open system and a closed system using a real example.
  \item Trace a single carbon atom through at least three reservoirs of the carbon cycle.
  \item Evaluate a short case study for evidence of feedback loops.
\end{enumerate}
\end{objectivesbox}

\begin{center}
\begin{minipage}{0.47\textwidth}
\begin{infobox}{\faClock\ \ Estimated Time}
Reading: 40 min\\
Activity: 20 min\\
Quiz: 10 min\\
\textbf{Total: about 70 minutes}
\end{infobox}
\end{minipage}\hfill
\begin{minipage}{0.47\textwidth}
\begin{infobox}{\faBook\ \ Reading Assignment}
\emph{Systems of the Living Earth}, Chapter 1, pp.\ 3--21, plus the recorded lecture
``Spheres in Motion'' (18 min) in the course shell.
\end{infobox}
\end{minipage}
\end{center}

\section{The Four Spheres}

Environmental science begins with a simple but powerful idea: the planet behaves as one
interconnected system rather than a collection of separate parts. Scientists divide this system
into four overlapping spheres. The \textbf{atmosphere} is the thin envelope of gases surrounding
the planet, extending roughly 100 km upward before thinning into space. The \textbf{hydrosphere}
includes every reservoir of water -- oceans, rivers, groundwater, ice, and atmospheric vapor. The
\textbf{lithosphere} is the solid outer shell of rock and soil, and the \textbf{biosphere} is the
sum of every living thing and the environments they inhabit.

No sphere operates in isolation. When a wildfire burns through a forest (biosphere), it releases
carbon into the atmosphere, alters the soil chemistry of the lithosphere, and can change local
rainfall patterns in the hydrosphere for years afterward. Recognizing these cross-sphere
interactions is the foundational skill this course builds toward across all eight modules.

\begin{center}
\begin{tikzpicture}[
  every node/.style={font=\small},
  sphere/.style={circle,draw=deepnavy,line width=0.9pt,fill=softgray,minimum size=1.4cm,align=center,font=\footnotesize\bfseries\color{deepnavy}}
]
  \node[sphere] (atmo) at (0,1.9) {Atmosphere};
  \node[sphere] (bio)  at (2.0,0) {Biosphere};
  \node[sphere] (litho) at (0,-1.9) {Lithosphere};
  \node[sphere] (hydro) at (-2.0,0) {Hydrosphere};

  \draw[->,tealaccent,line width=1pt] (atmo) -- (bio);
  \draw[->,tealaccent,line width=1pt] (bio) -- (litho);
  \draw[->,tealaccent,line width=1pt] (litho) -- (hydro);
  \draw[->,tealaccent,line width=1pt] (hydro) -- (atmo);

  \node[warmgold,font=\footnotesize\itshape] at (0,0) {constant exchange};
\end{tikzpicture}
\end{center}

\section{Open Systems, Closed Systems, and Feedback}

A \textbf{closed system} exchanges energy with its surroundings but not matter; Earth as a whole
is very nearly closed in this sense -- sunlight enters, heat radiates out, but almost no material
leaves or arrives. An \textbf{open system}, such as a single watershed, exchanges both energy and
matter freely with what surrounds it: rain enters, rivers carry sediment out, and organisms move
across its boundary daily.

Feedback loops determine whether a system resists or amplifies change. A \textbf{negative
feedback loop} dampens change and pushes a system back toward equilibrium -- for instance,
increased cloud cover reflecting sunlight and cooling a warming surface. A \textbf{positive
feedback loop} amplifies an initial change, such as melting sea ice exposing darker ocean water
that absorbs more heat and melts still more ice.

\begin{termsbox}
\begin{description}[leftmargin=3.6cm,itemsep=1pt,topsep=1pt]
  \item[Biosphere] All living organisms and the environments that sustain them.
  \item[Closed system] A system exchanging energy, but not matter, with its surroundings.
  \item[Feedback loop] A process where a system's output influences its own future behavior.
  \item[Reservoir] A location where a substance such as carbon is stored for a period of time.
  \item[Residence time] The average duration a unit of matter remains in a given reservoir.
\end{description}
\end{termsbox}

\section{Case Study: The Carbon Cycle}

Carbon moves continuously among the atmosphere, oceans, soil, and living tissue. A carbon atom
released by a factory smokestack may be absorbed by ocean surface water within a decade, taken up
by phytoplankton, sink to the deep ocean as organic debris, and remain there as sediment for
thousands of years before geological uplift and erosion return it to the surface. Each reservoir
has a different \textit{residence time} -- the atmosphere holds carbon for roughly 5 years on
average, while deep ocean sediments can hold it for millennia.

\begin{activitybox}
Choose one everyday object within arm's reach that contains carbon (a wooden pencil, a cotton
shirt, a plastic pen). In 3--4 sentences, trace the likely path that carbon took through at least
two spheres before becoming part of that object, and identify one reservoir where it may have
had a long residence time. Post your answer to the Module 1 discussion board before attempting
the quiz below.
\end{activitybox}

\begin{quizbox}
\begin{enumerate}[label=\arabic*.]
  \item Which sphere would sea ice, groundwater, and atmospheric humidity all belong to?
  \begin{enumerate}[label=(\alph*)]
    \item Lithosphere \quad (b) Hydrosphere \quad (c) Biosphere \quad (d) Atmosphere
  \end{enumerate}
  \item A single isolated lake that receives rainfall and loses water only to evaporation, with
    no matter otherwise entering or leaving, is best described as:
  \begin{enumerate}[label=(\alph*)]
    \item An open system \quad (b) A closed system \quad (c) A feedback loop \quad (d) A reservoir
  \end{enumerate}
  \item Melting Arctic ice exposing darker water that absorbs more solar energy, causing further
    melting, is an example of:
  \begin{enumerate}[label=(\alph*)]
    \item Negative feedback \quad (b) A closed system \quad (c) Positive feedback \quad (d) Residence time
  \end{enumerate}
  \item List two spheres a carbon atom might pass through between being emitted by a wildfire and
    being stored in deep ocean sediment.
\end{enumerate}
\end{quizbox}

% ================= CHAPTER 2 =================
\chapter{Populations, Resources, and Human Impact}

\begin{objectivesbox}
By the end of this chapter, you will be able to:
\begin{enumerate}[label=\arabic*.]
  \item Distinguish between renewable and non-renewable resources with an example of each.
  \item Explain carrying capacity and predict what happens when a population exceeds it.
  \item Compare two competing strategies for managing a shared resource.
\end{enumerate}
\end{objectivesbox}

\begin{center}
\begin{minipage}{0.47\textwidth}
\begin{infobox}{\faClock\ \ Estimated Time}
Reading: 45 min\\
Activity: 25 min\\
Quiz: 10 min\\
\textbf{Total: about 80 minutes}
\end{infobox}
\end{minipage}\hfill
\begin{minipage}{0.47\textwidth}
\begin{infobox}{\faBook\ \ Reading Assignment}
\emph{Systems of the Living Earth}, Chapter 2, pp.\ 23--44, plus the case brief
``The Apex Fishery Collapse'' posted in the course shell.
\end{infobox}
\end{minipage}
\end{center}

\section{Resources and Carrying Capacity}

Resources are commonly classified as \textbf{renewable} -- those that regenerate on human
timescales, such as timber or fresh groundwater recharge -- or \textbf{non-renewable}, those that
form over geological timescales far slower than we consume them, such as coal or phosphate rock. A
resource can behave as renewable in principle yet be depleted in practice if extraction outpaces
regeneration, as happened to the once-abundant Apex Bank fishery in the case brief for this
module. Every population, human or otherwise, exists within a \textbf{carrying capacity}: the
maximum size an environment can sustain indefinitely given available food, water, and space.
Populations that briefly exceed it typically experience a sharp correction -- famine, disease, or
emigration -- before stabilizing near the original limit, unless resource availability itself changes.

\begin{termsbox}
\begin{description}[leftmargin=3.6cm,itemsep=1pt,topsep=1pt]
  \item[Carrying capacity] The maximum population an environment can sustain long-term.
  \item[Ecological footprint] The land and resource area required to sustain a given lifestyle.
  \item[Non-renewable resource] A resource that regenerates far slower than it is consumed.
  \item[Tragedy of the commons] Overuse of a shared resource when no single party bears the full cost.
\end{description}
\end{termsbox}

\begin{activitybox}
Compare two solutions to a shared-resource problem of your choosing (a community well, a public
fishery, shared grazing land). Note whether each relies on regulation, pricing, or community
norms, and predict which better prevents a tragedy of the commons in that case.
\end{activitybox}

\begin{quizbox}
\begin{enumerate}[label=\arabic*.]
  \item A fishery harvested faster than fish populations can reproduce is best classified as:
  \begin{enumerate}[label=(\alph*)]
    \item A non-renewable resource in practice \quad (b) Always non-renewable \quad (c) A closed system
  \end{enumerate}
  \item When a population briefly exceeds carrying capacity, the most likely outcome is:
  \begin{enumerate}[label=(\alph*)]
    \item Permanent stabilization above the limit \quad (b) A correction back toward the limit
  \end{enumerate}
  \item Name one factor that increases an individual's ecological footprint.
\end{enumerate}
\end{quizbox}

% ================= ANSWER KEY =================
\chapter*{Answer Key}
\addcontentsline{toc}{chapter}{Answer Key}

\begin{infobox}{\faCheck\ \ Chapter 1 -- Self-Check Answers}
\begin{enumerate}[label=\arabic*.]
  \item (b) Hydrosphere
  \item (b) A closed system
  \item (c) Positive feedback
  \item Sample answer: biosphere to atmosphere (combustion releases carbon dioxide), then
    atmosphere to hydrosphere (dissolves into ocean surface water) before settling into
    lithosphere sediment.
\end{enumerate}
\end{infobox}

\vspace{10pt}

\begin{infobox}{\faCheck\ \ Chapter 2 -- Self-Check Answers}
\begin{enumerate}[label=\arabic*.]
  \item (a) A non-renewable resource in practice
  \item (b) A correction back toward the limit
  \item Sample answers: higher meat consumption, larger dwelling size, more frequent air travel,
    or greater per-capita energy use.
\end{enumerate}
\end{infobox}

\vspace{16pt}
\noindent{\small\color{inkgray}\faEnvelope\ \ Questions about scoring or content should go to
your instructor via the course messaging tool, not personal email, so responses count toward
your participation record. \faPhone\ \ Technical support: Meridian IT Helpdesk, available
Monday--Friday during posted office hours.}

\end{document}
